A repair neighborhood is defined relative to a feasible incumbent
$(\bar x,\bar y)$. Let $U\subseteq R$ contain the optional locations whose
visit decisions may change, and let $F\subseteq E$ contain the routing edges
whose decisions may change. The repair problem consists of the reference
formulation with the additional restrictions
\begin{align}
 x_e&=\bar x_e &&e\in E\setminus F,\label{eq:repair-fixed-edges}\\
 y_i&=\bar y_i &&i\in V\setminus U.\label{eq:repair-fixed-visits}
\end{align}
The incumbent satisfies these restrictions and therefore provides a feasible
fallback. Optional locations outside $U$ keep their incumbent visit status;
mandatory locations remain selected through the reference formulation.
The choice of $F$ must be stated together with $U$: releasing a visit alone
need not permit its removal when the required reconnection edges remain fixed.

This neighborhood separates the feasibility of the mathematical problem from
the behavior of a time-limited implementation. The incumbent remains
mathematically feasible even when a solver does not return a new solution.
An implementation can preserve it explicitly and validate every returned cycle
against the original visits, coverage sets, prize quota, and edge costs.
A returned feasible solution gives a valid PCCTP upper bound. In contrast,
a lower bound for the restricted repair problem need not be a lower bound
for the unrestricted PCCTP, because the restrictions may exclude better tours.

\subsection{Contraction of retained paths}
Let $H$ be the graph formed by the incumbent edges fixed to one in
\eqref{eq:repair-fixed-edges}. Since these edges belong to a simple cycle,
each nontrivial component of $H$ is a path, unless the entire incumbent cycle
is retained. Every retained path must occur in every feasible repair. Its
interior vertices already have degree two, while each endpoint requires one
additional incident edge. An isolated selected vertex instead requires two
incident edges. An isolated released vertex may be selected or omitted.

The path components can consequently be represented by contracted nodes with
their endpoint incidences retained explicitly. All routing edges fixed to zero
remain unavailable. A retained internal edge is not made available again merely
because both of its endpoints appear in a boundary set. The cost of fixed
edges is collected in the constant
\begin{equation}
 c_{\mathrm{fixed}}=\sum_{e\in E\setminus F:\,\bar x_e=1}c_e.
\end{equation}
Prizes and demand coverage of retained path vertices are accounted for in the
corresponding mandatory component. Degree requirements at the original endpoints
and connectivity between selected components complete the contracted model.
The endpoint representation is needed because an ordinary degree-two constraint
on a contracted node would not distinguish which original endpoint an edge uses.

The integer correspondence follows by contracting the fixed paths of a feasible
repair tour. Conversely, a connected contracted solution satisfying the endpoint
requirements expands uniquely along each retained path, up to its traversal
direction. The expanded routing edges satisfy the original degree equations and
all fixed decisions. Preserving the original prize and coverage contributions
then yields a feasible repair tour with objective equal to the contracted edge
cost plus $c_{\mathrm{fixed}}$. If a retained component is already the incumbent
cycle, its vertices have no remaining degree capacity. Connectivity precludes
adding a separate cycle, so the incumbent routing cycle is the only feasible
routing solution of that repair problem.

The contraction is used as a representation of a specified neighborhood, rather
than as a claim that the neighborhood is new. Its correctness does not imply
that a chosen neighborhood contains an improving solution, that a short solve
will find one, or that contraction improves overall running time. Those are
separate empirical questions. In particular, a short offline repair from an
incumbent generated earlier is distinct from an end-to-end algorithm whose
construction, selection, model assembly, and repair all share one time budget.
